\chapter{软件工程基础与需求分析}

\paragraph*{学习目标}

学完本章，读者应能够：

\begin{enumerate}
\item 说明软件生命周期各阶段的目标、输入和交付物；
\item 根据需求稳定性、风险、合规和反馈频率选择合适的生命周期模型；
\item 组织干系方识别、需求获取、分类、优先级排序与验证；
\item 使用系统流程图、数据流图、数据字典和决策表描述水利业务；
\item 编制结构完整、可追踪、可验证的软件需求规格说明。
\end{enumerate}

\paragraph*{引言}

需求分析的目的不是尽可能多地记录用户意见，而是形成各方能够共同确认、设计人员能够实现、测试人员能够验证的需求基线。水利项目还涉及跨部门数据、现场设备、专业模型和高风险控制，更需要明确对象、单位、时间、权限、异常和验收证据。本章以“清源灌区用水计划申报与配水”为主要示例，逐步完成从访谈到需求规格说明的工作。

\section{软件生命周期}

软件生命周期是软件从提出构想到停止使用的全过程\cite{ISO12207,IEEE1074,sommerville2015,bourque2014,zhang2013}。生命周期活动可以迭代或并行，但其责任与交付物不能省略。

\subsection{需求与规划阶段}

团队识别业务问题、干系方、范围、约束和成功标准，评估技术、经济、进度、安全与运维可行性。核心交付物包括范围说明、干系方清单、业务场景、需求基线和初步验收标准。需求分析的具体方法见2.3与2.4节。

\subsection{架构与详细设计阶段}

团队把需求转为系统边界、模块职责、数据模型、接口、权限和部署方案，并记录质量属性权衡。设计不仅覆盖正常路径，还要定义网络中断、数据异常、模型失败和无权限操作等异常路径。

\subsection{实现、集成与验证阶段}

实现阶段依据接口契约编码和配置；集成阶段按可运行切片连接前端、后端、数据、模型与设备；验证阶段检查产品是否正确实现规格，确认阶段检查产品是否解决真实业务问题。单元、接口、系统、性能、安全和用户验收测试分别提供不同证据。

\subsection{部署、运行与维护阶段}

上线前准备环境、数据迁移、培训、回滚和应急方案；运行中监控服务、数据、模型和业务闭环。维护包括缺陷修复、适应环境变化、功能增强与预防性改进。每次变更都要评估影响、更新文档并回归测试。

\subsection{退役与数据处置阶段}

系统停止使用时，需要迁移业务、归档或销毁数据、撤销账号和密钥、保留审计证据，并确认下游系统不再依赖旧接口。退役是受控活动，不是简单关闭服务器。表\ref{tab:sdlc-deliverables}把上述各阶段并排列出：第二列是该阶段必须回答的问题，第三列是回答之后应当留下的交付物。判断一个阶段是否真正完成，看的是第三列有没有产出可供下一阶段引用的东西，而不是日程表上的时间是否走完。

\begin{table}[htbp]
\centering
\caption{软件生命周期阶段、关键问题与交付物}
\label{tab:sdlc-deliverables}
\begin{tabular}{p{2.5cm}p{5.0cm}p{5.1cm}}
\toprule
阶段 & 关键问题 & 典型交付物 \\
\midrule
需求与规划 & 为什么做、为谁做、边界是什么 & 范围、场景、需求基线、验收标准 \\
架构与设计 & 如何分解、如何协作、如何失败 & 架构图、数据模型、接口与异常流 \\
实现与集成 & 组件是否按契约工作 & 代码、配置、构建产物、集成记录 \\
验证与确认 & 是否符合规格并解决业务问题 & 测试用例、缺陷、验收记录 \\
部署与运行 & 如何上线、监控、恢复和演化 & 发布包、运行手册、监控与工单 \\
退役 & 如何迁移依赖并处置数据 & 迁移、归档、销毁和关闭证明 \\
\bottomrule
\end{tabular}
\end{table}

\section{软件生命周期模型与选型}

生命周期模型规定活动如何排序、反馈和控制。模型不是固定模板；选型应服务于风险与协作，而不是追逐名称。

\subsection{瀑布模型}

瀑布模型按需求、设计、实现、验证和运行依次推进，阶段评审和文档边界清楚。它适用于需求相对稳定、合规审查严格、软硬件接口明确的部分，但反馈较晚，早期误解可能在后期放大。图\ref{fig:waterfall-model}中那条从“验证”折回“设计”的虚线，正是为缓解这一点而加的评审反馈；但它只跨越一个阶段，如果误解发生在需求环节，回退距离仍然很长。

\begin{figure}[H]
\centering
\begin{tikzpicture}[node distance=6mm and 7mm,
  model/.style={draw=Primary,rounded corners,fill=PrimaryLight,minimum width=2.0cm,
    minimum height=.75cm,align=center}]
\node[model] (r) {需求};
\node[model,right=of r] (d) {设计};
\node[model,right=of d] (b) {实现};
\node[model,right=of b] (t) {验证};
\node[model,right=of t] (o) {运行};
\draw[->,draw=Primary,thick] (r)--(d);
\draw[->,draw=Primary,thick] (d)--(b);
\draw[->,draw=Primary,thick] (b)--(t);
\draw[->,draw=Primary,thick] (t)--(o);
\draw[->,draw=Accent,dashed,thick,bend right=35] (t) to node[below,font=\scriptsize]{问题反馈} (d);
\end{tikzpicture}
\caption{带评审反馈的瀑布模型}
\label{fig:waterfall-model}
\end{figure}

\subsection{原型模型}

原型用于澄清用户难以直接表达的交互、报表或空间展示需求。团队快速构造可评估版本，收集反馈后修订需求。原型必须注明是探索性还是演化性；探索性原型不应在缺少架构、安全和测试的情况下直接上线。图\ref{fig:prototype-model}把这一过程画成小循环：构造、演示、收集反馈、修订需求，若干轮之后才进入正式开发。循环次数与退出条件应当事先约定，否则原型容易变成没有验收标准的长期返工。

\begin{figure}[H]
\centering
\begin{tikzpicture}[node distance=6mm and 8mm,
  model/.style={draw=Primary,rounded corners,fill=PrimaryLight,minimum width=2.2cm,
    minimum height=.75cm,align=center}]
\node[model] (ask) {初步需求};
\node[model,right=of ask] (quick) {快速设计};
\node[model,right=of quick] (proto) {构建原型};
\node[model,right=of proto] (review) {用户评价};
\node[model,below=of proto] (baseline) {形成需求基线};
\draw[->,draw=Primary,thick] (ask)--(quick);
\draw[->,draw=Primary,thick] (quick)--(proto);
\draw[->,draw=Primary,thick] (proto)--(review);
\draw[->,draw=Accent,thick,bend left=28] (review) to node[above,font=\scriptsize]{反馈} (quick);
\draw[->,draw=Primary,thick] (review)--(baseline);
\end{tikzpicture}
\caption{原型模型以用户反馈澄清需求}
\label{fig:prototype-model}
\end{figure}

\subsection{迭代与增量模型}

增量模型把能力分成若干可运行交付，例如先交付测点查询，再交付质量标记、预警和工单。迭代模型允许同一能力反复改进。两者常结合使用：每个增量经过需求、设计、实现和验证，小步获得真实反馈。图\ref{fig:incremental-model}把这种安排画了出来：每个增量都完整走过需求到验证，交付的是能独立运行、能独立验收的能力，而不是把做了一半的功能提前上线。

\begin{figure}[H]
\centering
\begin{tikzpicture}[node distance=5mm,
  inc/.style={draw=Primary,rounded corners,fill=PrimaryLight,minimum width=10cm,
    minimum height=.7cm,align=left}]
\node[inc] (i1) {增量1：测点、观测查询与质量码};
\node[inc,below=of i1] (i2) {增量2：规则预警、确认与工单};
\node[inc,below=of i2] (i3) {增量3：模型预演、会商与报告};
\draw[->,draw=Primary,thick] (i1)--(i2);
\draw[->,draw=Primary,thick] (i2)--(i3);
\end{tikzpicture}
\caption{智慧水利平台的可运行增量示例}
\label{fig:incremental-model}
\end{figure}

\subsection{螺旋模型}

螺旋模型以风险为中心，每轮识别目标和备选方案、分析风险、开发验证并规划下一轮。它适用于模型不确定、设备集成复杂或安全风险较高的大型项目，但需要成熟的风险分析和项目治理能力。

\subsection{敏捷与迭代交付}

敏捷强调通过短周期、可运行软件和持续反馈适应变化\cite{agile2001}。业务用户和客户代表应持续参与评审，并对需求优先级和验收结果拥有决策权；跨职能交付团队据此拆分增量、演示成果并记录反馈。运维人员可提供运行约束和故障数据，但敏捷的成立取决于持续反馈与可决策的协作机制。水利项目采用敏捷时，仍需保留法规、标准、架构、安全和验收所要求的文档与审批。

\subsection{模型选择}

选型不是挑一个名字，而是先判断风险集中在哪里，再确认团队能否承担该模型要求的治理成本。表\ref{tab:lifecycle-comparison}把前面几种模型放在需求稳定性、优势与风险三个维度下并列比较：螺旋模型要求成熟的风险分析能力，敏捷要求客户代表持续在场并能拍板，这些都是前提条件而非可选项，缺了前提，模型本身并不会自动生效。

\begin{table}[htbp]
\centering
\caption{生命周期模型选择对照}
\label{tab:lifecycle-comparison}
\begin{tabular}{p{2.0cm}p{3.4cm}p{3.8cm}p{3.4cm}}
\toprule
模型 & 适合情形 & 主要优势 & 主要风险 \\
\midrule
瀑布 & 需求稳定、审查严格 & 阶段和文档清楚 & 反馈较晚 \\
原型 & 交互或展示需求模糊 & 快速澄清理解 & 原型被误作产品 \\
增量/迭代 & 能力可拆分、需要早交付 & 小步验证、降低集成风险 & 整体架构被局部需求侵蚀 \\
螺旋 & 风险高、方案不确定 & 显式管理风险 & 成本与治理要求高 \\
敏捷 & 需求变化快、用户可持续参与 & 反馈频繁、适应变化 & 合规和架构治理被忽视 \\
\bottomrule
\end{tabular}
\end{table}

清源灌区项目可采用“总体基线+增量交付”：先按合同和行业约束确认范围、数据与安全基线，再按两到三周迭代交付用水计划、审批、配水和执行反馈。设备控制接口单独设置严格评审和联调窗口。

\section{需求获取与理解}

需求获取发生在深入分析之前。团队先理解业务语言、责任和当前流程，再判断哪些意见需要转化为系统需求。

\subsection{干系方与业务场景}

清源灌区的关键干系方包括用水户、管理站、调度人员、审批人、计量维护人员、平台运维人员和外部数据提供方。每类角色拥有不同目标和权限；“用户”不能作为一个笼统角色。

业务场景使用“触发—角色—前置条件—主流程—异常流—结果”描述。例如：

\begin{quote}
灌溉季开始前，用水户提交分时用水计划；管理站检查地块、作物、需水量和渠道能力。资料不全时退回补充；通过后进入调度审批。批准计划生成配水任务，执行流量与计划偏差超过容限时触发复核。
\end{quote}

\subsection{需求获取方法}

访谈适合了解职责与例外；现场观察适合发现实际操作与制度文本的差异；工作坊适合跨部门澄清术语与边界；问卷适合大范围收集偏好；文档分析适合提取规程、表单、接口和历史问题；原型适合验证交互与可视化。团队应组合方法并记录来源。

每条原始意见保留提出人、时间、业务理由和关联资料。例如，不能只写“系统应支持快速查询”，而应追问查询对象、时间范围、并发人数、可接受响应时间、数据完整性和失败提示。

\subsection{需求分类与优先级}

需求可分为业务需求、用户需求、功能需求、质量属性、数据与接口需求以及合规约束。分类帮助发现遗漏，但不能代替优先级决策。

MoSCoW方法把需求分为Must、Should、Could和本次Won't。Must表示没有它就无法验收、无法合规或无法安全运行；Should具有高价值但存在可接受替代；Could在时间允许时实现；Won't表示明确不进入当前范围。优先级必须记录理由、依赖和批准人，并在迭代或变更时复核。表\ref{tab:moscow-example}按这四档整理了清源灌区的一组需求。读表时应重点看“理由”一列：同样标为Must，“缺了它无法通过合规验收”与“用户呼声很高”是完全不同的依据，只有前者在范围压缩时不可让步。

\begin{table}[htbp]
\centering
\caption{清源灌区需求优先级示例}
\label{tab:moscow-example}
\begin{tabular}{p{2.0cm}p{6.6cm}p{4.0cm}}
\toprule
级别 & 需求 & 理由 \\
\midrule
Must & 批准计划前校验渠道能力与可供水量 & 防止生成不可执行任务 \\
Must & 所有审批和调整保留人员、时间与理由 & 责任和审计要求 \\
Should & 地图联动显示地块与配水状态 & 提升研判效率，可由列表降级 \\
Could & 语音输入查询条件 & 非核心便利功能 \\
Won't & 本期自动下发闸门控制 & 控制安全论证不在当前范围 \\
\bottomrule
\end{tabular}
\end{table}

\section{需求分析的任务与原则}

需求分析把获取到的材料转为一致、完整、可行、可验证的需求集合。

\subsection{核心任务}

\textbf{明确边界。} 定义系统内外部角色、设备、外部平台和人工活动，避免把组织管理问题全部推给软件。

\textbf{建立对象和数据语义。} 统一地块、渠道、闸门、用水计划、配水任务和执行记录的身份、关系、单位与时间。

\textbf{描述功能与异常流。} 除正常申报审批外，还要覆盖重复提交、资料缺失、可供水不足、设备离线、计划过期和撤回。

\textbf{量化质量属性。} 把“及时”“稳定”“安全”转为响应时间、可用率、恢复目标、权限规则和审计留存等可验证指标。

\textbf{定义接口与约束。} 说明外部数据来源、协议、频率、责任边界、错误处理和版本策略。

\textbf{建立可追踪关系。} 需求连接到来源、设计、接口、测试和验收证据，变更时能够评估影响。

\subsection{基本原则}

需求应当必要、无歧义、一致、可行、可验证、可追踪并具有优先级\cite{wiegers2013,robertson2012}。一条需求尽量表达一个可识别义务，使用“应”描述必须满足的行为，避免“可能包括”“尽可能”“用户友好”等模糊措辞。

需求还要区分事实、假设与设计决定。例如“灌区现有闸门控制接口只读”是现状约束，“网络可用率不低于某值”是待确认假设，“采用PostgreSQL”是架构决定，三者不能都伪装成业务需求。

\subsection{需求验证与冲突处理}

团队通过评审、原型、模型走查和验收用例验证需求。冲突由业务价值、安全、合规、成本和依赖共同决策，而不是由声音最大的人决定。变更后更新需求版本、影响分析和批准记录。表\ref{tab:requirement-quality-check}把这些要求整理成逐条可勾选的检查项，用于评审前自查。它的作用是把“需求写得好不好”这种主观判断，换成“这一条能否被测试用例验证、是否只有唯一解释、能否追溯到提出人与来源”这类当场就能回答的问题。

\begin{table}[htbp]
\centering
\caption{需求质量检查清单}
\label{tab:requirement-quality-check}
\begin{tabular}{p{2.4cm}p{7.2cm}p{3.0cm}}
\toprule
属性 & 检查问题 & 证据 \\
\midrule
必要 & 不实现会损害哪个业务目标或约束？ & 业务理由 \\
无歧义 & 角色、对象、条件、动作和结果是否唯一？ & 场景走查 \\
一致 & 是否与术语、数据、接口或其他需求冲突？ & 冲突记录 \\
可行 & 技术、成本、进度和现场条件是否允许？ & 可行性验证 \\
可验证 & 能否设计明确的通过/失败判据？ & 验收用例 \\
可追踪 & 能否定位来源、设计和测试？ & 追踪矩阵 \\
\bottomrule
\end{tabular}
\end{table}

\section{结构化分析}

结构化分析从系统边界出发，自顶向下分解处理、数据流、数据存储和外部实体。它适合说明业务信息如何流动，不用于表达所有时间、控制和组织关系。复杂控制状态可在后续设计中使用状态机等方法补充。

\subsection{系统流程图}

系统流程图关注业务步骤、判断、输入输出和责任衔接。图\ref{fig:order-process-flowchart}描述用水计划从申报到执行反馈的闭环：资料不完整时从“完整性检查”退回“提交计划”；审批不通过时返回“修订计划”；审批通过后才生成执行任务并接收回执。这样的顺序让每个节点都有明确责任人与可审计状态。

绘制流程图之前先统一符号，否则同一个菱形在不同人笔下可能表示条件判断，也可能表示人工审核环节，图就失去了作为交流依据的价值。表\ref{tab:system-flow-symbols}给出本书采用的符号约定。

\begin{longtable}{p{2.5cm}p{3.2cm}p{6.6cm}}
\caption{系统流程图常用符号}\label{tab:system-flow-symbols}\\
\toprule
符号 & 名称 & 用途 \\
\midrule
\endfirsthead
\toprule
符号 & 名称 & 用途 \\
\midrule
\endhead
圆角矩形 & 开始/结束 & 表示流程入口、正常结束或受控终止 \\
矩形 & 处理 & 表示申报、校验、审批、生成任务等活动 \\
菱形 & 判断 & 表示资料是否完整、方案是否批准等条件 \\
平行四边形 & 输入/输出 & 表示表单输入、报告输出或外部消息 \\
箭头 & 流向 & 表示业务执行顺序；回退箭头应标明条件 \\
\bottomrule
\end{longtable}

按上述符号约定绘制的用水计划流程见图\ref{fig:order-process-flowchart}。图中两条回退箭头都标注了触发条件，这正是流程图区别于示意图的地方：回退路径必须说清在什么条件下发生、退回到哪个节点，否则读图的人无法判断异常情况该由谁接手。

\begin{figure}[H]
\centering
\begin{tikzpicture}[node distance=6mm and 12mm,
  proc/.style={draw=Primary,rounded corners,fill=PrimaryLight,minimum width=2.5cm,
    minimum height=.75cm,align=center},
  decision/.style={diamond,draw=Accent,fill=Accent!10,aspect=2.0,align=center,inner sep=1pt}]
\node[proc] (submit) {提交用水计划};
\node[decision,below=of submit] (complete) {资料完整？};
\node[proc,below=of complete] (check) {校核需水与渠道能力};
\node[decision,below=of check] (approve) {批准？};
\node[proc,below=of approve] (task) {生成配水任务};
\node[proc,below=of task] (execute) {执行与计量回传};
\node[proc,right=of complete] (supplement) {补充资料};
\node[proc,right=of approve] (revise) {修订计划};
\node[proc,left=of execute] (review) {偏差复核};
\draw[->,draw=Primary,thick] (submit)--(complete);
\draw[->,draw=Primary,thick] (complete)--node[left,font=\scriptsize]{是} (check);
\draw[->,draw=Primary,thick] (complete)--node[above,font=\scriptsize]{否} (supplement);
\draw[->,draw=Accent,thick] (supplement) |- (submit);
\draw[->,draw=Primary,thick] (check)--(approve);
\draw[->,draw=Primary,thick] (approve)--node[left,font=\scriptsize]{是} (task);
\draw[->,draw=Primary,thick] (approve)--node[above,font=\scriptsize]{否} (revise);
\draw[->,draw=Accent,thick] (revise) |- (check);
\draw[->,draw=Primary,thick] (task)--(execute);
\draw[->,draw=Primary,thick] (execute)--(review);
\draw[->,draw=Accent,dashed,thick,bend left=30] (review) to node[above,font=\scriptsize]{形成下一周期依据} (submit);
\end{tikzpicture}
\caption{用水计划申报、审批、配水与反馈流程}
\label{fig:order-process-flowchart}
\end{figure}

\subsection{数据流图}

数据流图（DFD）只表达数据如何在外部实体、处理和数据存储之间流动。上下文图把系统视为一个处理，先确定边界；再逐层分解，但父图输入输出必须与子图保持平衡。图\ref{fig:dfd-level-zero}是本例的上下文图：整个系统被画成单个处理，围绕它的外部实体确定了系统边界，后续任何一层分解都不能凭空引入图中没有出现的外部实体。

\begin{figure}[H]
\centering
\begin{tikzpicture}[node distance=12mm and 23mm,
  entity/.style={draw=Primary,fill=PrimaryLight,minimum width=2.3cm,minimum height=.8cm,align=center},
  process/.style={ellipse,draw=Accent,fill=Accent!10,minimum width=3.0cm,minimum height=1.2cm,align=center}]
\node[entity] (farmer) {用水户};
\node[process,right=of farmer] (system) {灌区用水计划\\与配水系统};
\node[entity,right=of system] (manager) {调度与审批人员};
\node[entity,below=of system] (meter) {计量与闸门系统};
\draw[->,draw=Primary,thick] (farmer)--node[above,font=\scriptsize]{申报/补充资料} (system);
\draw[->,draw=Primary,thick] (system)--node[below,font=\scriptsize]{状态/批准计划} (farmer);
\draw[->,draw=Primary,thick] (manager)--node[above,font=\scriptsize]{审批/调整} (system);
\draw[->,draw=Primary,thick] (system)--node[below,font=\scriptsize]{待办/分析结果} (manager);
\draw[->,draw=Primary,thick] (system)--node[right,font=\scriptsize]{配水任务} (meter);
\draw[->,draw=Primary,thick] (meter)--node[left,font=\scriptsize]{执行回执} (system);
\end{tikzpicture}
\caption{灌区用水计划与配水系统上下文图}
\label{fig:dfd-level-zero}
\end{figure}

把上下文图中的单个处理展开，就得到图\ref{fig:dfd-level-one}的一级数据流图。检查分解是否正确的办法是做平衡核对：一级图中所有跨越系统边界的数据流，必须与上下文图中进出系统的数据流逐条对应，既不能多出一条，也不能漏掉一条。

\begin{figure}[H]
\centering
\begin{tikzpicture}[node distance=7mm and 6mm,
  process/.style={ellipse,draw=Primary,fill=PrimaryLight,minimum width=2.15cm,minimum height=.9cm,align=center,font=\small},
  store/.style={draw=Primary,double,minimum width=2.25cm,minimum height=.7cm,align=center,font=\small}]
\node[process] (receive) {1.0 接收申报};
\node[process,right=of receive] (validate) {2.0 校验与分析};
\node[process,right=of validate] (approve) {3.0 审批计划};
\node[process,right=of approve] (dispatch) {4.0 生成配水};
\node[store,below=of validate] (plan) {D1 用水计划};
\node[store,below=of approve] (resource) {D2 水源与工程能力};
\node[store,below=of dispatch] (receipt) {D3 执行记录};
\draw[->,draw=Primary,thick] (receive)--(validate);
\draw[->,draw=Primary,thick] (validate)--(approve);
\draw[->,draw=Primary,thick] (approve)--(dispatch);
\draw[<->,draw=Primary,thick] (validate)--(plan);
\draw[<->,draw=Primary,thick] (validate)--(resource);
\draw[->,draw=Primary,thick] (dispatch)--(receipt);
\draw[->,draw=Accent,dashed,thick] (receipt)--(validate);
\end{tikzpicture}
\caption{灌区用水计划与配水系统一级数据流图}
\label{fig:dfd-level-one}
\end{figure}

\subsection{数据字典}

数据字典定义数据流、数据存储和数据元素的结构与约束。每个字段的说明栏应给出编码格式、数据类型、单位、来源、取值域和校验规则，例如计划编号采用项目编码加年度序号，状态字段只能取受控枚举，流量字段注明精度与上限。清晰的字段契约能让接口、数据库和测试用例使用同一口径。表\ref{tab:water-plan-dictionary}按这个要求整理了用水计划的主要字段。填写时最容易省略的是“取值域”与“校验规则”两栏，而恰恰是这两栏决定了接口能否拒绝非法输入——把它们留空，等于把校验责任推给了下游每一个调用方。

\begin{longtable}{p{2.9cm}p{2.0cm}p{1.7cm}p{4.8cm}}
\caption{用水计划数据字典示例}\label{tab:water-plan-dictionary}\\
\toprule
数据项 & 类型/长度 & 单位 & 说明与校验 \\
\midrule
\endfirsthead
\toprule
数据项 & 类型/长度 & 单位 & 说明与校验 \\
\midrule
\endhead
planId & 字符串/26 & — & 全局稳定ID，不使用名称替代 \\
applicantId & 字符串/18 & — & 关联用水户主数据，必须有效 \\
canalId & 字符串/12 & — & 关联渠道对象编码 \\
startTime & 时间戳 & UTC+8 & 与endTime构成左闭右开时段 \\
endTime & 时间戳 & UTC+8 & 必须晚于startTime \\
{\small requestedVolume} & 十进制 & m$^3$ & 大于0，精度和上限由业务规则规定 \\
maxFlow & 十进制 & m$^3$/s & 不得超过渠道批准能力 \\
priority & 枚举 & — & emergency、must、normal等受控值 \\
status & 枚举 & — & draft、submitted、approved、rejected、executing、closed \\
version & 整数 & — & 每次修订递增，批准后不可覆盖 \\
\bottomrule
\end{longtable}

数据流也要定义。例如“执行回执”至少包含任务ID、设备ID、指令时间、实际启闭时间、目标流量、实测流量、质量码和回执状态。涉及水工设备动作时，接口、日志和界面统一使用“闸门启闭”，并分别记录目标开度、实际开度、启闭状态与操作指令，避免把动作命令和观测结果混为一项。

\subsection{决策表}

当业务结果由多个条件组合决定时，决策表比长段落更清楚。下表示例只用于需求澄清，真实阈值应由批准规则确定。表\ref{tab:water-plan-decision}给出三个条件的组合及其对应动作。检查决策表是否完备的办法是数组合：$n$个二值条件共有$2^n$种取值，表中各条规则覆盖的组合数，加上用“无关”合并掉的组合数，应当正好等于这个总数；对不上就说明有组合没有被任何规则接住。

\begin{table}[htbp]
\centering
\caption{用水计划审批决策表示例}
\label{tab:water-plan-decision}
\begin{tabular}{p{4.2cm}cccc}
\toprule
条件/动作 & 规则1 & 规则2 & 规则3 & 规则4 \\
\midrule
资料完整 & 是 & 是 & 否 & 是 \\
水量可满足 & 是 & 否 & — & 是 \\
渠道可用 & 是 & 是 & — & 否 \\
\midrule
批准计划 & $\checkmark$ & & & \\
进入方案调整 & & $\checkmark$ & & $\checkmark$ \\
退回补充资料 & & & $\checkmark$ & \\
\bottomrule
\end{tabular}
\end{table}

\section{软件需求规格说明与验收}

GB/T 9385-2008《计算机软件需求规格说明规范》目前继续有效，可作为组织软件需求规格说明（SRS）的依据\cite{gbt93852008}。课程项目可采用下列精简目录，但不能省略关键内容：

\begin{enumerate}
\item 引言：目的、范围、术语、参考资料；
\item 总体描述：产品背景、用户特征、运行环境、假设与约束；
\item 外部接口：用户、软件、硬件和通信接口；
\item 具体需求：功能、数据、质量属性、安全与合规；
\item 业务规则与异常处理；
\item 验收标准与需求追踪关系；
\item 附录：数据字典、模型、待确认问题与变更记录。
\end{enumerate}

一条可验证需求可以写为：“当已批准用水计划进入执行时段且现场计量有效时，系统应每5分钟比较计划流量与实测流量；连续3个周期超过批准容限时，创建唯一偏差事件并通知值班员。” 其中角色、条件、频率、判断、结果和去重要求都可以设计测试。

\begin{table}[htbp]
\centering
\caption{需求追踪矩阵示例}
\label{tab:requirements-traceability}
\begin{tabular}{p{2.0cm}p{3.6cm}p{3.2cm}p{3.2cm}}
\toprule
需求ID & 来源/业务目标 & 设计或接口 & 验收用例 \\
\midrule
REQ-PLAN-01 & 用水户申报场景 & POST /water-plans & TC-PLAN-01～04 \\
REQ-PLAN-07 & 渠道能力约束 & 计划校验服务 & TC-CAP-01～03 \\
REQ-AUDIT-02 & 审批责任追踪 & 审计事件接口 & TC-AUDIT-05 \\
REQ-DEV-03 & 执行偏差复核 & 偏差规则与工单 & TC-DEV-01～06 \\
\bottomrule
\end{tabular}
\end{table}

评审SRS时，应让业务、开发、测试、运维和安全角色从各自角度走查。评审通过形成基线；后续变更需提交原因、影响、优先级和批准记录，而不是直接修改共享文档。

\paragraph{本章图表导读}
表\ref{tab:sdlc-deliverables}按生命周期阶段列出输入、关键问题和交付物。
图\ref{fig:waterfall-model}、图\ref{fig:prototype-model}和图\ref{fig:incremental-model}分别帮助比较计划驱动、反馈澄清和可运行增量。
表\ref{tab:lifecycle-comparison}用于按风险与反馈频率选型；表\ref{tab:moscow-example}展示需求优先级。
表\ref{tab:requirement-quality-check}把完整性、可验证性和一致性转为检查项；表\ref{tab:system-flow-symbols}解释流程图符号。
图\ref{fig:order-process-flowchart}呈现用水计划闭环；图\ref{fig:dfd-level-zero}和图\ref{fig:dfd-level-one}分别刻画上下文与一级数据流。
表\ref{tab:water-plan-dictionary}给出字段契约；表\ref{tab:water-plan-decision}说明条件组合；表\ref{tab:requirements-traceability}把需求连接到设计和验收证据。

\section*{小结}

软件生命周期覆盖需求、设计、实现、验证、运行、维护与退役；瀑布、原型、增量、螺旋和敏捷模型适合不同风险与反馈条件。需求工作应先获取和理解，再分类、排序、分析与验证。结构化分析用流程图、DFD、数据字典和决策表建立一致模型，最终形成可追踪、可验证的SRS。

\section*{章末交付物}

以“清源灌区用水计划申报与配水”为对象，提交：干系方清单、3个业务场景、MoSCoW优先级表、上下文DFD、一级DFD、数据字典、决策表、精简SRS和需求追踪矩阵。

\section*{思考题与练习题}

\begin{enumerate}
\item 软件生命周期为什么包含退役与数据处置？
\item 对需求稳定、接口明确且审查严格的设备接入项目，生命周期模型应如何选择？说明理由与风险。
\item 探索性原型为什么不能直接视为生产系统？
\item 为灌区用水计划系统识别6类干系方，并分别写出一个核心目标。
\item 使用MoSCoW方法排列5条需求，写明优先级理由和依赖。
\item 将“系统应快速显示渠道状态”改写为可验证需求。
\item 按图\ref{fig:order-process-flowchart}增加“计划过期”异常路径，说明回退位置。
\item 检查图\ref{fig:dfd-level-zero}与图\ref{fig:dfd-level-one}是否保持输入输出平衡。
\item 使用DFD描述“用水计划申报—审批—配水—执行反馈”，不得把施工进度等非数据对象画成数据流。
\item 为“执行回执”补充不少于8个数据字典字段，标明类型、单位和校验规则。
\item 编写一张至少含3个条件、4条规则的配水审批决策表。
\item 按2.6节目录编制一份5～8页精简SRS，并为5条关键需求建立追踪矩阵。
\end{enumerate}
